\section{The frequency cutoff}\label{cc:app:cutoff}
Fix a radial $\psi\in C_c^\infty(\R^3)$ with $0\le\psi\le1$, $\psi=1$ on
$|\zeta|\le1$ and $\supp\psi\subset\{|\zeta|\le2\}$. For $\Lambda\ge1$ let
$S_\Lambda$ multiply the spatial Fourier mode $e^{2\pi i\xi\cdot x}$,
$\xi\in\mathbb Z^3$, by $\psi(\xi/\Lambda)$.
Subsection~\ref{cc:uf:trunc:sec:truncation} takes $u=S_\Lambda V$ as the
next parent. The frequencies of $u$ are at most $2\Lambda$, and those of
$u\cdot\nabla u$ are at most $4\Lambda$. Thus \eqref{cc:uf:trunc:eq:lp}
bounds derivatives of both $u$ and its residual, while
\eqref{cc:uf:trunc:eq:tail} controls the error $u-V$.

\begin{lemma}\label{cc:uf:trunc:lem:cutoff}
$S_\Lambda$ commutes with derivatives, preserves parity and divergence
freeness, and maps into fields supported on frequencies at most
$2\Lambda$. Its operator norms on $L^\infty$ and on $\Holder{0}$ are bounded
in terms of $\psi$ alone. For every integer $m\ge0$ and every
normed-space-valued $f$ supported on frequencies at most $4\Lambda$,
\begin{equation}\label{cc:uf:trunc:eq:lp}
 \sum_{|\beta|=m}\norm{\partial^\beta f}_\infty\le C_m\Lambda^m\norm f_\infty .
\end{equation}
For all integers $0\le s\le M$ and every $f\in\Holder{M}$,
\begin{equation}\label{cc:uf:trunc:eq:tail}
 \norm{(\Id-S_\Lambda)f}_{\Holder{s}}
 \le C_M\Lambda^{s-M}\norm f_{\Holder{M}}.
\end{equation}
The constants $C_m$ and $C_M$ depend only on their index and on $\psi$.
\end{lemma}

\begin{proof}
The symbol is scalar, even and independent of time, so the cutoff
preserves parity and divergence freeness and commutes with differentiation.
Its support gives the frequency bound. For the norm estimates we use
the convolution kernel as in the Euclidean Bernstein inequalities
\cite[Appendix~A, (A.2)--(A.6)]{TaoDispersive}. Let $K=\check\psi$,
a Schwartz function, and set $K_\Lambda(y)=\Lambda^3K(\Lambda y)$.
To compute $S_\Lambda f$,
convolve the periodic extension of $f$ to $\R^3$ with $K_\Lambda$.
Since $\norm{K_\Lambda}_{L^1}=\norm K_{L^1}$, applying convolution
to $f$ and to its spatial increments gives the two operator bounds.
If $f$ has frequencies at most $4\Lambda$, then $f=S_{4\Lambda}f$.
The identity $\norm{\partial^\beta K_{4\Lambda}}_{L^1}
=(4\Lambda)^{|\beta|}\norm{\partial^\beta K}_{L^1}$ gives
\eqref{cc:uf:trunc:eq:lp} also for normed-space-valued $f$.

For \eqref{cc:uf:trunc:eq:tail} put $k=M-s$. When $k=0$, apply the
$\Holder{0}$ bound to the derivatives of $f$ through order $s$.
For $k\ge1$, the kernel has integral one and all its positive-order
moments vanish because $\psi$ is constant near zero. Thus in
$f-S_\Lambda f$ only the remainder after expanding $f(x-y)$ at $y=0$
through degree $k-1$ contributes. Its $\Holder{s}$ norm in $x$ is bounded
by $C_M|y|^k\norm f_{\Holder{M}}$, by Taylor's formula for each derivative
through order $s$ and its spatial increments. Integrating this bound
against $|K_\Lambda(y)|$ gives \eqref{cc:uf:trunc:eq:tail}, since
$\int_{\R^3}|y|^k|K_\Lambda(y)|\dd y
=\Lambda^{-k}\int_{\R^3}|y|^k|K(y)|\dd y$.
\end{proof}
